% paper/construction_capacity.tex -- the rigid-group rule's evidence base:
% 2024 capacity signals for the sectors the programme presses on.
%
% Companion documents: paper/framing_gamma.tex (Section sec:capacity argues
% the capacity hypothesis; this note supplies the external evidence that
% turns that argument into a testable rule), docs/decisions/ADR-0022
% (the sectoral-labour-market closure), docs/grouping_rule_evidence.md (the
% full evidence table with the 2026-09-20 source verification),
% docs/decisions/ADR-0024 (the programme total G0 = 40,300 context).
%
% Evidence vintage: 2024, matching the programme vintage of the executed
% generation (impulses.csv 2024 row, G0 = 40,300 EUR m at 2019 prices).
%
% Version 1 (September 2026). No LaTeX in the analysis container: compile
% on the Mac (pdflatex; no .bib dependency).

\documentclass[11pt]{article}

\usepackage[english]{babel}
\usepackage[a4paper,top=2.2cm,bottom=2.2cm,left=3cm,right=3cm]{geometry}
\usepackage{amsmath,amssymb,amsthm}
\usepackage{booktabs}
\usepackage{xcolor}
\providecolor{revisionV1}{HTML}{800080}
\providecolor{revisionV2}{HTML}{CC5500}
\providecolor{revisionV3}{HTML}{006400}
\providecolor{revisionV4}{HTML}{0000CD}
\providecolor{revisionV5}{HTML}{808080}
\providecolor{revisionV6}{HTML}{8B0000}
\usepackage[colorlinks=true,allcolors=blue]{hyperref}

\setlength\parindent{0pt}
\setlength{\parskip}{0.45em}

\newcommand{\Lbar}{\bar{L}}

\title{Construction capacity and the rigid-group rule\\\\[0.4em]
\large The 2024 evidence base for which sectors are rigid in the
sectoral-labour-market closure}
\author{Hermes Agent (Lt.\ Cmdr.\ Data), for Prof.\ Dr.\ J.\ Kapeller}
\date{September 20, 2026}

\begin{document}
\maketitle

\begin{center}\small
\textbf{Version 1} (September 2026)
\end{center}

\begin{abstract}
The executed sectoral family carries one discretionary modelling choice:
\textit{which} sectors' labour supplies are treated as rigid
($\eta_{s,i} = 0$). The choice flips the sign of the programme's employment
effect --- rigid programme sectors: $-0.34\,\%$; rigid largest half:
$+0.02\,\%$; uniform: $+0.13\,\%$ at F2 --- so it must be a rule fixed before
the runs, not a partition read off the results. The companion framing note
motivates the rule by capacity pressure in the sectors the programme presses
on; this note supplies the external, pre-run evidence for that motivation:
four 2024 indicators --- the vacancy survey, the skilled-worker-shortage
barometer, order-backlog data and the establishment panel --- triangulate
\textit{how hard it is to get workers quickly} in the programme's sectors,
and the rule they imply selects, from data alone, the seven programme
sectors. The limits are stated as carefully as the numbers: the indicators
select the rigid \textit{set}, they do not estimate the elasticity; their
sector coverage is coarser than the model's; and 2024 was a downturn, so the
claim rests on elasticity signals, not on utilisation.
\end{abstract}

\section{The choice at stake}\label{sec:stake}

The sectoral closure (ADR-0022) replaces one aggregate real-wage equation by
$N$ sectoral markets,
\begin{equation}
L^{cm}_i(p, y_i, w_i) \;=\; \bar{L}_i\,
\Big[ \frac{w_i / \Pi(p)}{\bar{w}_i / \bar{\Pi}} \Big]^{\eta_{s,i}},
\end{equation}
and the executed equipment ladder in the flow table shows what the family
delivers: at $\eta_{s,i} = 0$ for every $i$ the system is the $\eta = 0$
endpoint (Proposition 3 of \texttt{paper/equivalence.tex}), the rigid corner
of the family. The level of the elasticity is unidentified under
demand-only shocks --- the ladder is a sensitivity band --- but the
\textit{membership} of the rigid set is identified up to the rule that
selects it, and the selection is consequential: making the seven programme
sectors rigid ($\eta_{s,i} = 0$ there) reverses the employment effect
relative to a uniform elasticity and to rigidity on the largest half by
baseline employment (framing note, Table~1 of the rigidity sweep).

ADR-0022 therefore demands ``a rule \ldots{} preferred to a hand-picked
list, to forestall the tuning objection''. The companion note's answer is
the capacity-pressure hypothesis: the programme is a demand push aimed at
the economy's project-capacity chain, and the short-run margin in the
sectors it presses on is capacity --- skilled trades cannot be trained and
project capacity cannot be built within the horizon of the impulse. That
hypothesis is \emph{testable} against external sectoral evidence
(vacancies, overtime, backlogs in the construction trades), which is what
this note provides. The rule has two conditions, both computable before any
model run:

\begin{quote}\small
\textbf{The rule.} A sector's labour supply is rigid ($\eta_{s,i} = 0$)
iff (i) its programme incidence share is at least 1\,\% of the programme
total \emph{and} (ii) its skilled-worker-shortage signal reaches the
goods-producing average or above.
\end{quote}

\section{The programme's incidence (condition i)}\label{sec:incidence}

From the 2024 row of the impulse table (\texttt{cbase2/data_raw/impulses.csv},
G0 = 40{,}300 EUR m): the programme lives in seven of the seventy-one
sectors --- 63.4\,\% of it in specialised construction works and the rest in
its material suppliers.

\begin{center}
\begin{tabular}{lc}
\toprule
Sector (VGR) & 2024 impulse share \\
\midrule
Specialised construction works & 63.4\,\% \\
Rubber and plastics products & 9.4\,\% \\
Ceramic products, processed stone and clay & 8.0\,\% \\
Chemicals and chemical products & 7.6\,\% \\
Glass and glassware & 6.1\,\% \\
Machinery & 4.3\,\% \\
Electrical equipment & 1.3\,\% \\
\bottomrule
\end{tabular}
\end{center}

Condition~(i) selects exactly these seven sectors; the next sector in the
ranking has zero incidence. The composition is stable over the 2024--2050
horizon (max $\lvert \Delta \psi \rvert = 0.092$ against the 2024 vector,
worst year 2050, when machinery's share rises to 13.5\,\%).

\section{The capacity signals (condition ii), 2024}\label{sec:signals}

Four indicators, three national series and one pair of regional panels.
None of them measures a labour-supply elasticity; together they triangulate
how hard it is to get workers quickly in the programme's sectors --- the
short-run margin the rigidity assumption is about.

\begin{center}
\begin{tabular}{llll}
\toprule
Indicator & Value (2024) & Series & Coverage \\
\midrule
Skilled-worker-shortage share, & 28.9\,\% (Q4, & KfW-ifo-Fachkr\"afte- & WZ 41.2, 42.x, \\
building main trades & survey Oct.) & barometer (Dec. 2024) & 43.1, 43.9 \\
Skilled-worker-shortage share, & 20.6\,\% (Q4) & same & manufacturing \\
manufacturing (aggregate) & & & broad \\
Skilled-worker-shortage share, & 33.1\,\% (Q4) & same & rubber/plastics \\
rubber and plastics & & & (WZ 22) \\
Open positions, construction & 108{,}000 (Q4) & IAB-Stellenerhebung, & construction \\
& & IAB-Monitor 4/2024 & sector \\
Vakanzrate (national) & 3.2\,\% (Q4) & IAB-Stellenerhebung & all sectors \\
Order backlog, building main & 3.7 months (Jun.) & ifo Konjunktur- & Bauhauptgewerbe \\
trades & & umfragen (2024) & branches \\
Working-time accounts, & 20\,\% $\to$ 62\,\% & IAB-Betriebspanel & Hesse; West \\
construction (2023 $\to$ 2024) & & Report Hessen 2024 & 37 $\to$ 58\,\% \\
Hiring persistence, construction & 54\,\% of firms; 72\,\% & IAB-Betriebspanel & Brandenburg \\
(H1 2024) & of those unfilled & 2024 & \\
\bottomrule
\end{tabular}
\end{center}

\textbf{Vacancy survey.} The IAB's quarterly establishment survey counts
all positions firms are trying to fill, including those outside the
employment agency's register. Q4/2024: 1.40 million open positions
nationally, of which the construction sector holds 108{,}000 --- a top-five
sector stock. The national \emph{Vakanzrate} (immediately open positions per
100 workers demanded) stood at 3.2\,\%. The reading is a supply symptom:
even after a year in which vacancies fell 19\,\% (the 2024 downturn),
construction firms keep posting jobs and keep failing to fill them.

\textbf{Skilled-worker-shortage barometer.} The KfW-ifo barometer asks about
9{,}000 firms per quarter one question --- is your business currently
restricted by a lack of skilled workers? --- and the reported figure is the
share of firms answering yes. Q4/2024: 28.9\,\% in the building main trades
against 20.6\,\% in manufacturing broadly and 33.1\,\% in rubber and
plastics (one of the six material sectors, and the most affected
manufacturing branch in the report). Construction runs at roughly 1.4 times
the manufacturing rate at a moment when manufacturing itself is slack. The
measure is self-reported current restriction and is cyclical (it fell from a
37.0\,\% peak in Q4/2022).

\textbf{Order backlogs.} The ifo survey reports how many months of work
order books cover at current capacity; 3.7 months (June 2024, Tiefbau 4.1)
means that if new orders stopped, firms still have close to four months of
scheduled work. This indicator is included as the \emph{counter-signal}: it
fell from 4.8 months in early 2022, so the 2024 story is not one of bursting
order books --- the rigidity claim rests on the labour-elasticity signals,
not on utilisation.

\textbf{Establishment panel.} Two regional 2024 reports (Hesse, Brandenburg)
serve as directional evidence, not national estimates. (i) Working-time
accounts in construction jumped from 20\,\% to 62\,\% of firms over
2023--2024 (37\,\% $\to$ 58\,\% in the West): when order books are volatile,
firms keep headcount and buffer through \emph{hours} --- hours are the
flexible margin, employment the rigid one, which is exactly the vertical
supply curve the closure assumes. (ii) In Brandenburg 54\,\% of construction
firms had a hiring need in H1/2024 (third-highest of all sectors) and 72\,\%
of those left positions unfilled: recruitment failure persists even in a
weak year.

\section{The rule, and its discipline}\label{sec:discipline}

The rigidity in the model is a property of the \emph{slope} of a sector's
labour supply --- an elasticity that none of these series measures. The
rule therefore uses the indicators only to \emph{select the set}; the value
of $\eta_s$ itself comes from the supply arm (ADR-0023), which identifies it
under a technology shock. Four qualifications, stated for the record:

\begin{itemize}
\item \textbf{Data gaps.} The 2024 sub-branch shortage shares for
  machinery and electrical equipment are not published in the barometer we
  retrieved; under the rule, the defensible pre-run default is the
  five-sector core \{specialised construction, plastics, ceramics,
  chemicals, glass\} (94.5\,\% of programme incidence), with a weaker
  threshold variant (the full seven) as the sensitivity cell.
\item \textbf{Mapping.} The model's sectors are VGR (Destatis 2019
  input-output classification); the reports use WZ/NACE groupings.
  ``Specialised construction works'' maps to WZ 43, while the
  Bauhauptgewerbe indicator covers WZ 41.2, 42.x, 43.1, 43.9 --- a superset.
\item \textbf{Vintage and cycle.} All signals are 2024, the programme
  vintage; the downturn caveat is the backlog row, and the rule's thresholds
  are set on elasticity-type signals, not utilisation.
\item \textbf{Scope.} As in ADR-0022, a welfare difference across rules is
  not a statement about wage policy; the rule only answers which sectors'
  supply curves are treated as inelastic.
\end{itemize}

\section{What this delivers}\label{sec:delivers}

For the preregistration, the rule and its thresholds are fixed from public
2024 statistics \emph{before} any model outcome is computed; the map from
data to the rigid set is a calculation, not a choice. The same evidence
ranges are the natural probability measure for the sectoral Sobol
decomposition of the family's response (IDEA-0003), so the evidence base
serves both the rule and the attribution step. The full table with sources
and verification dates lives in \texttt{docs/grouping\_rule\_evidence.md}
(v1, 2026-09-20).

\section*{Sources}

\begin{itemize}
\item KfW Research / ifo Institut, \emph{KfW-ifo-Fachkr\"aftebarometer}
  December 2024 (survey October 2024).
\item IAB, \emph{IAB-Stellenerhebung} (quarterly); IAB-Monitor
  Arbeitskr\"aftebedarf 4/2024 (IAB-Forum, 2025).
\item ifo Institut, \emph{Konjunkturumfragen Bauhauptgewerbe},
  Auftragsbestand in Monaten, 2024.
\item IAB-Betriebspanel 2024: Report Hessen 2024 (HNA/iwak); Betriebspanel
  Brandenburg 2024.
\end{itemize}

\end{document}